\documentclass[10pt]{article}

\usepackage{fancyhdr}
\pagestyle{fancy}
\fancyhf{}
\fancyhead[C]{Working Note: Claims Not Verified}
\fancyfoot[C]{\thepage}
\renewcommand{\headrulewidth}{0.4pt}

\usepackage[margin=2cm]{geometry}
\usepackage[T1]{fontenc}
\usepackage[utf8]{inputenc}
\usepackage{amsmath,amssymb,amsthm,mathtools,mathrsfs}
\usepackage{tikz-cd}
\usepackage{tikz}
\usetikzlibrary{positioning,arrows.meta}
\usepackage{booktabs}
\usepackage{enumitem}
\usepackage{tabularx}
\usepackage{xcolor}
\usepackage[colorlinks=true,linkcolor=blue!55!black,citecolor=blue!55!black,urlcolor=blue!55!black]{hyperref}
\usepackage[round]{natbib}

\setlist{nosep,leftmargin=1.3em}
\setlength{\bibsep}{1pt}
\renewcommand{\bibfont}{\footnotesize}

\theoremstyle{plain}
\newtheorem{proposition}{Proposition}

\newcommand{\Hmu}[1]{H_\mu\!\left[#1\right]}
\newcommand{\occ}{d^{\pi}}
\newcommand{\E}{\mathbb E}
\newcommand{\Scal}{\mathcal S}
\newcommand{\Acal}{\mathcal A}
\newcommand{\Zcal}{\mathcal Z}
\newcommand{\Occ}{\mathscr D}
\newcommand{\Kb}{\bar K}
\newcommand{\Pb}{\bar P}
\newcommand{\id}{\mathrm{id}}
\newcommand{\copym}[1]{\mathrm{copy}_{#1}}
\newcommand{\pair}[2]{\langle #1,#2\rangle}
\newcommand{\Hom}{\operatorname{Hom}}
\newcommand{\open}[1]{\textcolor{red!70!black}{[Open: #1]}}

\title{The Learner as a Regulator:\\ Inference for Regulation in the Good-Regulator Lens}
\author{Working note, 7 October 2026}
\date{}

\begin{document}
\maketitle
\thispagestyle{fancy}

\noindent
Conant and Ashby say which regulator is good and nothing about how it is found.
This note poses the finding as a regulation problem of the same shape, so that their theorem applies a second time.
Everything is written as morphisms in $\mathbf{FinStoch}$, in the notation of Neil's \texttt{good\_regulator\_categorical.tex}: a state on $X$ is a morphism $I \to X$, $\copym{X}$ duplicates a wire, and $\pair{f}{g}$ pairs two morphisms on a copied input.
Section~\ref{sec:chain} first walks from their diagram to the named sequential problems, one arrow at a time.
The categorical notation from Section~\ref{sec:loops} on is there to see three things at once: which object the loop runs on, where the occupancy lives, and which wires the regulator reads as opposed to which wires the world uses.

\section{From Conant and Ashby to the sequential problems}
\label{sec:chain}

Start from Neil's deterministic form of the theorem: sets $D, S, A, Z$, a map $\varphi\colon D \to S$ from disturbance to state, a map $\psi\colon S \times A \to Z$ from state and action to outcome, a law $p$ on $D$, and a regulator that reads $S$ under error control, $\rho\colon S \to A$, or reads $D$ under cause control, $\rho\colon D \to A$.
The criterion is $\min H(Z)$.
The regulator is already a policy; what the moves below change is its domain, and the theorem says at every step that the regulator's model is the sufficient statistic of that domain.
Each move changes one arrow of the diagram and nothing else, Figure~\ref{fig:moves}.

\begin{figure}[t]
\centering
\tikzset{v/.style={inner sep=1pt, minimum size=5mm}, e/.style={-{Stealth[length=1.6mm]}}, c/.style={-{Stealth[length=1.6mm]}, dashed, gray}}
\begin{minipage}{0.48\textwidth}\centering
\begin{tikzpicture}
  \foreach \t/\x in {t/0, t+1/2.2} {
    \node[v] (D\t) at (\x,1.2) {$D_{\t}$};
    \node[v] (S\t) at (\x,0.2) {$S_{\t}$};
    \node[v] (A\t) at (\x,-0.8) {$A_{\t}$};
    \node[v] (Z\t) at (\x+1.1,-1.7) {$Z_{\t}$};
    \draw[e] (D\t) -- node[left, font=\tiny] {$\varphi$} (S\t);
    \draw[e] (S\t) -- (A\t);
    \draw[e] (S\t) -- (Z\t);
    \draw[e] (A\t) -- (Z\t);
    \draw[c] (D\t) to[bend right=55] (A\t);
  }
\end{tikzpicture}\\[1mm]
{\small (a) Move 1: occasions.}
\end{minipage}\hfill
\begin{minipage}{0.48\textwidth}\centering
\begin{tikzpicture}
  \foreach \t/\x in {t/0, t+1/2.2} {
    \node[v] (D\t) at (\x,1.2) {$D_{\t}$};
    \node[v] (S\t) at (\x,0.2) {$S_{\t}$};
    \node[v] (A\t) at (\x,-0.8) {$A_{\t}$};
    \node[v] (Z\t) at (\x+1.1,-1.7) {$Z_{\t}$};
    \draw[e] (S\t) -- (A\t);
    \draw[e] (S\t) -- (Z\t);
    \draw[e] (A\t) -- (Z\t);
    \draw[c] (D\t) to[bend right=55] (A\t);
  }
  \draw[e, thick] (Dt) -- node[above, font=\tiny] {$f$} (St+1);
  \draw[e, thick] (St) -- (St+1);
  \draw[e, thick] (At) -- (St+1);
\end{tikzpicture}\\[1mm]
{\small (b) Move 2: $\varphi$ becomes $f$.}
\end{minipage}\\[3mm]
\begin{minipage}{0.48\textwidth}\centering
\begin{tikzpicture}
  \foreach \t/\x in {t/0, t+1/2.2} {
    \node[v] (D\t) at (\x,1.2) {$D_{\t}$};
    \node[v] (S\t) at (\x,0.2) {$S_{\t}$};
    \node[v] (A\t) at (\x,-0.8) {$A_{\t}$};
    \node[v] (Z\t) at (\x+1.1,-1.7) {$Z_{\t}$};
    \draw[e] (S\t) -- (A\t);
    \draw[e] (S\t) -- (Z\t);
    \draw[e] (A\t) -- (Z\t);
    \draw[c] (D\t) to[bend right=55] (A\t);
  }
  \draw[e] (Dt) -- (St+1);
  \draw[e] (St) -- (St+1);
  \draw[e] (At) -- (St+1);
  \draw[e, thick] (Dt) -- node[above, font=\tiny] {$g$} (Dt+1);
\end{tikzpicture}\\[1mm]
{\small (c) Move 3: $D$ has memory.}
\end{minipage}\hfill
\begin{minipage}{0.48\textwidth}\centering
\begin{tikzpicture}
  \foreach \t/\x in {t/0, t+1/2.2} {
    \node[v] (D\t) at (\x,1.2) {$\theta_{\t}$};
    \node[v] (S\t) at (\x,0.2) {$S_{\t}$};
    \node[v] (A\t) at (\x,-0.8) {$A_{\t}$};
    \node[v] (Z\t) at (\x+1.1,-1.7) {$Z_{\t}$};
    \draw[e] (S\t) -- (A\t);
    \draw[e] (S\t) -- (Z\t);
    \draw[e] (A\t) -- (Z\t);
    \draw[e, thick] (D\t) to[bend left=35] (Z\t);
    \draw[c] (D\t) to[bend right=55] (A\t);
  }
  \draw[e] (Dt) -- (St+1);
  \draw[e] (St) -- (St+1);
  \draw[e] (At) -- (St+1);
  \draw[e, thick] (Dt) -- node[above, font=\tiny] {$\delta_\eta$} (Dt+1);
\end{tikzpicture}\\[1mm]
{\small (d) Move 4: $D = \theta$.}
\end{minipage}
\caption{Four moves from Conant and Ashby's diagram, each adding the thick arrows and changing nothing else.
The dashed grey arrow is the cause-controlled reading, in which the regulator reads $D_t$ as well as $S_t$; without it the regulator is error-controlled.}
\label{fig:moves}
\end{figure}

\paragraph{Move 1: occasions are time steps.}
Conant and Ashby's $p$ is a frequency over the occasions on which the regulator acts.
Index them: $d_t \sim p$ independently, $s_t = \varphi(d_t)$, $a_t \sim \pi(\cdot \mid s_t)$, $z_t \sim K(\cdot \mid s_t,a_t)$, and nothing carries over from one occasion to the next.
This is a contextual bandit with context $s_t$, and the criterion is the entropy of the outcome frequency, $\min_\pi H\big(\sum_s p(s)\,\pi(a \mid s)\,K(z \mid s,a)\big)$.
The plain bandit is $\varphi$ constant: the disturbance then acts through $\psi$ alone, $z = \psi(d,a)$, and the error-controlled regulator is a single arm.
Under cause control the context is $d_t$ itself.

\paragraph{Move 2: the action moves the next state.}
A system that persists has a state at the next occasion that depends on the previous state and action.
Replace $\varphi\colon D \to S$ by $f\colon S \times A \times D \to S$ and change nothing else.
With $d_t$ independent, $P(s' \mid s,a) = \sum_d p(d)\, f(s' \mid s,a,d)$, the frequency over occasions is the occupancy, and the error-controlled regulator is a stationary policy on the MDP $(\Scal,\Acal,P)$ with outcome kernel $K$.
This is Section~3.3 of the paper.
The cause-controlled regulator reads $d_t$ as well, which is the MDP on $(s,d)$ with $d$ independent across steps.

\paragraph{Move 3: the disturbance has memory.}
Replace $d_t \sim p$ independent by $d_{t+1} \sim g(\cdot \mid d_t)$.
The pair $(s,d)$ is Markov, and the error-controlled regulator sees $s$ alone, so it faces a POMDP with hidden state $d$.
A memoryless policy on $s$ loses the polytope and may need to be stochastic \citep{Singh-learning94}.
The belief $b_t$ over $d_t$ restores the MDP on $(s,b)$, and that belief is the internal model of the disturbance generator \citep{Francis-internal76}.
Under cause control the regulator reads $d_t$ and the problem stays the MDP on $(s,d)$; no belief is needed.
Partial observation of $s$ itself, $o = h(s)$, is the same move with $s$ in place of $d$, and the belief over $s$ is Wentworth's case.

\paragraph{Move 4: the generator is the identity.}
Let $d_t = \theta$, constant, or drifting at rate $\eta$ by $\delta_\eta$.
The hidden disturbance is now the parameter of the kernels, $P_\theta$ and $K_\theta$, and the error-controlled regulator faces a POMDP whose hidden state never moves: the Bayes-adaptive MDP, with the posterior over $\theta$ as its belief \citep{Duff-optimal02}.
Under cause control $\theta$ is read and the problem is the known MDP of Move 2.
Without Move 2, so with $S$ trivial, it is the Bayesian bandit of \texttt{entropy\_budget\_empowerment.tex}.

\paragraph{What is and is not Conant and Ashby.}
Move 1 is their model: $p$ is already a frequency over occasions, and the contextual bandit is their diagram with the occasions written out.
Move 2 is the one extension of the diagram.
In their model the action reaches only $Z$; the arrow from $(S_t,A_t)$ into $S_{t+1}$ makes the regulator a cause of its own next state.
It stays inside their frame if the disturbance at occasion $t+1$ is taken to contain $(S_t,A_t)$, so that the regulator's own past is part of what disturbs it, and the theorem survives because the state marginal of the occupancy plays the role of $p(s)$.
Moves 3 and 4 do not change the diagram; they change what $S$ means.
Their $S$ is whatever the regulator reads, and under memory that is the history, so the theorem says the regulator is a function of the history's sufficient statistic, which is \citet{Wentworth-fixing21}.
Two slots are enlarged, $D$ by the regulator's past and $S$ by the history, and the outcome map, the criterion, concavity, determinism and the model as a quotient are untouched.
What does not survive is in the paper's assumption table: $\psi$ becomes the kernel $K$, $H(Z \mid S) = 0$ fails under noise, and "every" becomes "some" unless equivalence is taken on $(K,P)$.
Throughout, $\psi$ is the outcome map, as in Conant and Ashby and the paper; $\varphi$ is Neil's name for the arrow from $D$ to $S$, which the paper leaves unnamed.

\paragraph{Where the chain ends.}
Table~\ref{tab:cells} lists the cells.
In every one of them the criterion is the same, $\min_d \Hmu{dK}$ over the occupancy polytope of the policy's domain, the optimum is at a vertex, and the regulator is deterministic on that domain and a model of its sufficient statistic.
What the moves change is only the domain, from $s$ to $(s,d)$ to $(s,b)$, and so what the model is, from a quotient of $S$ to a belief over the generator to a posterior over the kernels.
Level~1 of Section~\ref{sec:loops} is the cell (Move 2, error); Level~2 is the cell (Move 4, error); the rest of this note is about those two cells.

\begin{table}[h]
\centering
\small
\begin{tabularx}{\textwidth}{@{}lXX@{}}
\toprule
Disturbance & Error-controlled: reads $s$ & Cause-controlled: reads $(s,d)$ \\
\midrule
Independent across steps (Moves 1--2) & Contextual bandit; with $f$, the MDP on $s$. Policy $\pi(a \mid s)$. Model: the coarsest quotient of $S$ that fixes $a$. & Contextual bandit with full context; with $f$, the MDP on $(s,d)$. Policy $\pi(a \mid s,d)$. Model: a quotient of $S \times D$. \\
Markov, generator $g$ (Move 3) & POMDP with hidden $d$; the belief MDP on $(s,b)$. Policy $\pi(a \mid s,b)$. Model: $b$, the internal model of $g$. & The MDP on $(s,d)$, as above. No belief. \\
Constant, or drifting by $\delta_\eta$ (Move 4) & Bayesian bandit; with $f$, the Bayes-adaptive MDP on $(s,b_\theta)$. Policy $\pi(a \mid s,b_\theta)$. Model: the posterior, the gooder regulator. & Known bandit; with $f$, the known MDP. \\
\bottomrule
\end{tabularx}
\caption{The sequential problems as cells of Conant and Ashby's diagram.
Rows are what the disturbance remembers; columns are what the regulator reads; "with $f$" is Move 2.
The objective and the vertex argument are the same in every cell.}
\label{tab:cells}
\end{table}

\section{Three loops}
\label{sec:loops}

\paragraph{Level 0: one decision.}
Neil's stochastic regulation problem is $(S, A, Z, \psi, p)$ with $p\colon I \to S$ exogenous and the closed loop
\begin{equation}
\begin{tikzcd}[column sep=large]
I \arrow[r,"p"] & S \arrow[r,"\copym{S}"] & S \otimes S \arrow[r,"\id \otimes \pi"] & S \otimes A \arrow[r,"K"] & Z
\end{tikzcd}
\qquad
\omega_\pi = K \circ (\id \otimes \pi) \circ \copym{S} \circ p .
\label{eq:level0}
\end{equation}
The map $\pi \mapsto \omega_\pi$ is affine, entropy is strictly concave, so some optimal $\pi$ is deterministic.
The regulator reads $S$ and nothing feeds back.

\paragraph{Level 1: the loop in time.}
Replace the outcome map by the pair $P\colon S \otimes A \to S$ and $K\colon S \otimes A \to Z$.
The arrow $P$ closes the diagram,
\begin{equation}
\begin{tikzcd}[column sep=large]
I \arrow[r,"p_\pi"] & S \arrow[r,"\copym{S}"] & S \otimes S \arrow[r,"\id \otimes \pi"] & S \otimes A \arrow[r,"K"] \arrow[ll, bend left=40, "P"] & Z
\end{tikzcd}
\label{eq:level1}
\end{equation}
and $p$ is no longer exogenous.
Write $T_\pi = P \circ (\id \otimes \pi) \circ \copym{S}\colon S \to S$ for one turn of the loop.
The state on the loop object is its fixed point, $p_\pi = T_\pi \circ p_\pi$ in the average case and $p_\pi = (1-\gamma)\rho_0 + \gamma\, T_\pi \circ p_\pi$ in the discounted case.
The occupancy is the state one wire later, $\occ = (\id \otimes \pi) \circ \copym{S} \circ p_\pi\colon I \to S \otimes A$, and the outcome law is $\omega_\pi = K \circ \occ$.
Two things changed against \eqref{eq:level0}.
The map $\pi \mapsto \occ$ is no longer affine, because $p_\pi$ depends on $\pi$ through the fixed point; the map $d \mapsto K \circ d$ still is, and the set of all $\occ$ is the polytope $\Occ$ with deterministic policies at its vertices \citep{Derman-finite70ch7}, so the vertex argument is run on $d$ instead of on $\pi$.
And the regulator still reads only $S$; the loop uses $P$ and $K$, but $\pi$ does not.

\paragraph{Level 2: the parameter joins the loop.}
Let $\theta \in \Theta$ index the kernels, $P_\theta$ and $K_\theta$, so that the world's step is a single morphism $\pair{P}{K}\colon \Theta \otimes S \otimes A \to S \otimes Z$ that reads the $\Theta$ wire.
The regulator cannot read $\Theta$.
It reads the history, and the history is summarised by the posterior $b_t = q(\theta \mid h_t)$, so the regulator is $\pi\colon S \otimes B \to A$ and carries a Bayes update $u\colon B \otimes S \otimes A \otimes S \otimes Z \to B$.
The loop object is $\Theta \otimes S \otimes B$ and one turn is
\begin{equation}
\begin{tikzcd}[column sep=large, row sep=large]
\Theta \otimes S \otimes B \arrow[r, "\id_\Theta \otimes \pair{\id}{\pi}"] & \Theta \otimes S \otimes B \otimes A \arrow[r, "\pair{\id}{\pair{P}{K}}"] & \Theta \otimes S \otimes B \otimes A \otimes S \otimes Z \arrow[d, "\mathrm{proj}_Z"] \arrow[ll, bend left=25, "\delta_\eta \otimes \mathrm{proj}_{S'} \otimes u"'] \\
& & Z
\end{tikzcd}
\label{eq:level2}
\end{equation}
where $\delta_\eta\colon \Theta \to \Theta$ is the drift, which keeps $\theta$ with probability $1-\eta$ and redraws it from $q_0$ otherwise, and $\mathrm{proj}_{S'}$ keeps the new state.
Call this turn $T_{\pi,\eta}$.
Its fixed point is a state on $\Theta \otimes S \otimes B$, the occupancy is that state one wire later, $d\colon I \to \Theta \otimes S \otimes B \otimes A$, and the outcome law is the $Z$ wire, $\omega = \mathrm{proj}_Z \circ \pair{\id}{\pair{P}{K}} \circ d$.
Setting $B = I$ and $\Theta = I$ recovers \eqref{eq:level1}, and setting $B = I$ with a point mass on $\Theta$ recovers it with $\theta$ known.
Table~\ref{tab:loops} records what each level reads.

\begin{table}[h]
\centering
\small
\begin{tabularx}{\textwidth}{@{}lXXXX@{}}
\toprule
Level & Loop object & Occupancy, a state on & Regulator reads & World's step reads \\
\midrule
0 & none, $p$ given & $S \otimes A$ & $S$ & $S \otimes A$ \\
1 & $S$ & $S \otimes A$ & $S$ & $S \otimes A$, through $P$ and $K$ \\
2 & $\Theta \otimes S \otimes B$ & $\Theta \otimes S \otimes B \otimes A$ & $S \otimes B$ & $\Theta \otimes S \otimes A$, through $P_\theta$ and $K_\theta$ \\
\bottomrule
\end{tabularx}
\caption{One occupancy per level, always the fixed-point state of the loop pushed one wire forward.
The regulator never reads $\Theta$; it reads $B$, and $B$ is its model of $\Theta$.}
\label{tab:loops}
\end{table}

\section{The regulator on the loop}
\label{sec:belief}

The question "does the outer level depend on $b$ and no longer on $(P,K)$" is answered by the wires in \eqref{eq:level2}.
The action depends on $(s,b)$ and on nothing else.
The outcome depends on $(\theta,s,a)$ and on nothing else.
The only route from $\theta$ to the action runs through past outcomes, into $h_t$, into $b_t$, which is Wentworth's diagram with $\theta$ as the system and $b$ as the memory \citep{Wentworth-fixing21}.
Figure~\ref{fig:wires} draws one step as a graph, which is the same information as \eqref{eq:level2} read wire by wire.

\begin{figure}[h]
\centering
\tikzset{v/.style={inner sep=1.5pt, minimum size=5mm}, e/.style={-{Stealth[length=1.8mm]}}}
\begin{tikzpicture}
  \node[v] (th) at (0,1.6) {$\theta_t$};
  \node[v] (th1) at (3.2,1.6) {$\theta_{t+1}$};
  \node[v] (s) at (0,0.4) {$s_t$};
  \node[v] (s1) at (3.2,0.4) {$s_{t+1}$};
  \node[v] (a) at (1.6,-0.6) {$a_t$};
  \node[v] (z) at (3.2,-0.6) {$z_t$};
  \node[v] (b) at (0,-1.8) {$b_t$};
  \node[v] (b1) at (3.2,-1.8) {$b_{t+1}$};
  \draw[e] (th) -- node[above, font=\scriptsize] {$\delta_\eta$} (th1);
  \draw[e] (th) -- (s1);
  \draw[e] (th) -- (z);
  \draw[e] (s) -- (s1);
  \draw[e] (s) -- (a);
  \draw[e] (b) -- (a);
  \draw[e] (a) -- (s1);
  \draw[e] (a) -- (z);
  \draw[e] (s) -- (z);
  \draw[e] (b) -- (b1);
  \draw[e] (z) -- (b1);
  \draw[e] (s1) -- (b1);
  \draw[e] (a) -- (b1);
\end{tikzpicture}
\caption{One turn of \eqref{eq:level2} as a graph.
Arrows into $a_t$ are what the regulator reads; arrows into $s_{t+1}$ and $z_t$ are what the world uses; arrows into $b_{t+1}$ are the Bayes update $u$.
The two levels are one regulator on $(s,b)$: the "outer" map $b \mapsto \pi(\cdot \mid \cdot, b)$ is $\pi$ curried, and the "inner" regulator is $\pi$ at a fixed $b$.}
\label{fig:wires}
\end{figure}

\begin{proposition}[The predictive law is the law of the run]
\label{prop:run}
Let $d$ be the occupancy of \eqref{eq:level2} and $d_{SBA}$ its marginal on $S \otimes B \otimes A$.
Then $d = (b \otimes \id) \circ d_{SBA}$, that is $d(\theta,s,b,a) = d_{SBA}(s,b,a)\, b(\theta)$, and the outcome law is $\omega = \Kb \circ d_{SBA}$ with $\Kb(z \mid s,b,a) = \E_{\theta \sim b} K_\theta(z \mid s,a)$.
\end{proposition}
\begin{proof}
$(s_t,b_t)$ is a function of $h_t$ and $b_t = q(\cdot \mid h_t)$ is the same on every history in the fibre, so $P(\theta \mid s_t,b_t) = b_t(\theta)$; multiply by $P(s_t,b_t)\,\pi(a \mid s_t,b_t)$, sum over $t$ with the occupancy weights, and push the $\Theta$ wire through $K_\theta$.
With the drift $\delta_\eta$ the posterior is over the current $\theta_t$ and the same argument applies.
\end{proof}

So the objective on this loop, $\min \Hmu{\Kb \circ d_{SBA}}$, is Conant and Ashby's criterion over everything the regulator does not control, $\theta$ included, with no approximation, and the learning phase is inside the occupancy because the chain on $(s,b)$ counts every step of the run.
The set of $d_{SBA}$ is again cut out by a flow constraint, so it is convex with deterministic policies at its extreme points, finite when $B$ is finite and countable for conjugate posteriors \citep{Altman-markov91}.

\begin{proposition}[The gooder regulator]
\label{prop:gooder}
Some optimal $\pi\colon S \otimes B \to A$ is deterministic, $H(A \mid S,B) = 0$, and in general no optimal regulator is deterministic in $S$ alone.
\end{proposition}

\noindent
This is Wentworth's theorem on the loop: the regulator's minimal model of $\theta$ is $b$, and its action is a function of the state and the model.
It also answers how the good regulator finds $\pi$.
It does not search.
The morphism $\pi\colon S \otimes B \to A$ is fixed, the belief moves by the fixed morphism $u$, and the policy on $S$ that an observer sees, $\pi(\cdot \mid \cdot, b_t)$, drifts with $b_t$.
Learning is the deterministic motion of the model wire inside a fixed regulator, which is the Bayes-adaptive MDP \citep{Duff-optimal02}.
Without the drift $\delta_\eta$ and with the average occupancy, the $B$ wire converges, the learning steps have weight zero, and the objective does not price them; the drift is the persistent disturbance of v1, applied to the $\Theta$ wire, and in the discounted case $\gamma$ supplies it as a reset with $1-\gamma = \eta$.

\section{Empirical Bayes: the search appears}
\label{sec:eb}

Computing the fixed $\pi$ of Proposition~\ref{prop:gooder} is intractable beyond bandits, so replace the $B$ wire by a point estimate $w$ and the Bayes update $u$ by maximum likelihood.
The inner step needs $\theta$ only through the outcome law, by the Gibbs identity
\begin{equation}
  \Hmu{p} = \min_q \E_{p}[-\ln q],
  \qquad
  \E_{\omega_\pi}[-\ln q] = -\sum_{s,a} \occ(s,a)\, r_q(s,a),
  \qquad
  r_q(s,a) = \E_{K(\cdot \mid s,a)}[\ln q(z)],
  \label{eq:gibbs}
\end{equation}
so $w$ parametrises $q$ and the regulator is $\pi(a \mid s,w)$ as in v1.
Now the loop of \eqref{eq:level1} runs with $\pi_k$ fixed, and a second loop runs on the hom-set,
\begin{equation}
\begin{tikzcd}[column sep=huge]
\pi_k \in \Hom(S,A) \arrow[r, "\text{run \eqref{eq:level1}}"] & d^{\pi_k} \arrow[r, "K"] & q_k = \omega_{\pi_k} \arrow[r, "\ln"] & r_{q_k} \arrow[r, "\text{solve MDP}"] & \pi_{k+1}
\end{tikzcd}
\label{eq:mm}
\end{equation}
The second arrow is the regulator producing its own data, the third is type-II maximum likelihood, since $q$ is the hyperparameter of the inner reward and is fit to the data the inner problem generates, and the last arrow is the search that Conant and Ashby never discuss.
In \eqref{eq:level2} both loops were one morphism $T_{\pi,\eta}$; here they are separated, the inner one run to its fixed point before the outer one moves.
At a fixed point $q^\star = \omega_{\pi^\star}$ and $\pi^\star$ solves control as inference with likelihood $q^\star$ \citep{Toussaint-probabilistic06,Levine-reinforcement18}.

\paragraph{What the separation costs.}
Each $\pi_k$ is scored by $\Hmu{\omega_{\pi_k}}$, the entropy of its own fixed point.
The run is the sequence, and its outcome law is, up to transients, $\bar\omega = \sum_k \lambda_k\, \omega_{\pi_k}$ with $\lambda_k$ the fraction of time under $\pi_k$, so by concavity
\begin{equation}
  \Hmu{\bar\omega} = \sum_k \lambda_k \Hmu{\omega_{\pi_k}} + I(Z;k).
  \label{eq:run-gap}
\end{equation}
MM descends the first term and does not see $I(Z;k)$, the information the outcome carries about which iterate produced it.
In \eqref{eq:level2} the iterate index is the $B$ wire, which is inside the loop object, so Proposition~\ref{prop:run} scores $\bar\omega$ and charges this term.
It is the across-episode term that separates the occupancy from SMiRL \citep{Berseth-smirl21,Mutti-importance22}, with the iterate in place of the episode.
\open{Whether $I(Z;k)$ is small for the MM loop in practice.}

\section{What must be modelled}
\label{sec:sf}

The outer regulator must model the second arrow of \eqref{eq:mm}, the map from a policy to its outcome law, and the minimal model of that map is not $\theta$.
With $\phi\colon Z \to \mathbb R^m$ the one-hot feature, or the feature family of v2, the map factors,
\begin{equation}
\begin{tikzcd}[column sep=huge]
\Hom(S,A) \arrow[r, "\pi \mapsto \psi^\pi"] \arrow[rr, bend right=20, "\pi \mapsto \omega_\pi"'] & \mathbb R^{S \times A \times m} \arrow[r, "(1-\gamma)\,\E_{\rho_0,\pi}"] & \mathrm{Dist}(Z)
\end{tikzcd}
\qquad
\psi^\pi = \E_K[\phi] + \gamma\, P^\pi \psi^\pi ,
\label{eq:sf}
\end{equation}
and the inner MDP's action value is $w^\top \psi^\pi$ up to the gauge term of v2.
Two parameters $\theta, \theta'$ with the same $\psi^\pi$ for every $\pi$ give the same bent arrow, so the coarsest quotient of $\Theta$ that preserves $\psi$ is the outer regulator's minimal non-faithful model.
Successor-feature abstractions are exactly the abstractions that predict every reward linear in $\phi$ \citep{Lehnert-successor19}, and $r_q$ in \eqref{eq:gibbs} is linear in $\phi$ for every $q$ in the family, so this is the quotient the paper's Section~4 should define: the outcome-bisimulation quotient of the state, relative to $Z$.
At level 2 the same reduction says the $B$ wire need only carry the posterior over $\psi$, not over $\theta$.
\open{Prove that the posterior over $\psi^\pi$ is sufficient for the level-2 objective; immediate for the bandit, where $\psi = \theta$.}

\section{The criterion as a fixed point, and what a model is}
\label{sec:fixedpoint}
Added 7 October 2026, after the discussion above.

\paragraph{A property of the argmin, not a method.}
Conant and Ashby state a property of optimal regulators, as Bellman optimality states a property of $V^\star$.
Neither says how to get there.
Bellman's statement is already a fixed-point equation; theirs is not, and the occupancy objective is what turns it into one.
Written out, the analogue of Bellman optimality is two coupled conditions,
\begin{equation}
  q^\star = d^{\pi^\star} K,
  \qquad
  \pi^\star \in \arg\max_\pi \sum_{s,a} \occ(s,a)\, \E_{K(\cdot \mid s,a)}\big[\ln q^\star(z)\big],
  \label{eq:fp}
\end{equation}
the second an ordinary Bellman problem with reward $r_{q^\star}$, the first closing it on the regulator's own outcome law.
A good regulator is a policy that is optimal for the reward defined by the outcomes it produces.
The method analogues follow: MM is policy iteration on the pair, the SARSA and natural-gradient actor of v1 is the Q-learning analogue, and the analogy breaks at uniqueness, since the objective is concave on the polytope and \eqref{eq:fp} has many solutions.
This is the fixed point that \texttt{frame.tex} calls variational EM and that Section~\ref{sec:eb} reaches from the empirical-Bayes side.

\paragraph{Model as the statistics of the fixed point.}
Under Bellman, model-free means carrying $Q$ and model-based means carrying $P$, both sufficient statistics for the fixed point.
Read the same way, "having a model" under the good-regulator criterion means carrying the statistics in which \eqref{eq:fp} is written, and the minimal model depends on what the regulator has to do with them.
\begin{itemize}
\item \emph{Act at the fixed point.} The reward $r_{q^\star}$ is fixed, so $Q$ for that reward, or the greedy action, suffices. This is Bellman's model-free level and Conant and Ashby's quotient of $S$.
\item \emph{Hold the fixed point.} The regulator also needs $q^\star$, to check that the outcomes it produces are the ones its reward was built from. $q$ is the pushforward $dK$, a histogram over $Z$ estimated from outcome samples; the occupancy itself is never needed. It is a model of the regulator's own outcome law, not of the environment.
\item \emph{Find the fixed point.} $q$ moves, so the reward moves, and $Q$ is needed for every $q$ along the way: either relearned per $q$, which is model-free on a moving reward and is v1's sampled reward $\ln \hat q(z_t)$, or carried as $\psi^\pi$, which gives $Q_w = w^\top \psi^\pi$ for every $q$ at once. $\psi$ is a model along the reward axis and model-free along the dynamics axis \citep{Lehnert-successor19}: it answers any reward linear in $\phi$ without new data and answers nothing about a different policy, which is why GPI is needed to move between policies \citep{Barreto-successor17}.
\end{itemize}
The environment model $P$ is never needed on any rung; what the top rung needs is $\psi$, a quotient of $P$ relative to $Z$, which is where "minimal, not faithful" gets its content here.
Our setting is the regime successor features were built for, a reward that moves at every step and dynamics that do not; the cost is the on-policy tie, $\psi$ is stale after a policy step until it tracks, which is why v1 keeps the actor rate small.
This ladder is the Part~3 fallback of \texttt{plan.md} (23 September), derived here from the criterion rather than from the algorithm.

\paragraph{Do we need occupancies, or can we use control as inference directly?}
Control as inference adds a binary optimality variable $O_t$ with $p(O_t = 1 \mid s,a) = r(s,a)$ \citep{Toussaint-probabilistic06,Levine-reinforcement18}.
A natural question is whether we can skip occupancies and minimise $H(O)$ instead.

\emph{A binary $O$ alone does not work.}
The long-run fraction of steps with $O = 1$ is the average reward $R(\pi)$, so the per-step entropy is $H(O) = h_b(R)$.
Over the achievable interval $R \in [R_{\min}, R^\star]$ this is minimised at whichever endpoint is further from $\tfrac12$.
In a sparse-reward task $R_{\min} = 0$ is achievable and $h_b(0) = 0$, so the minimiser is the policy that never succeeds.

\emph{Enlarging the alphabet of $O$ is exactly our reward channel.}
Let $O$ emit the good symbol with probability $r(s,a)$ and one of $N$ noise symbols uniformly otherwise.
Its per-step law under $\pi$ is the channel $K_{R(\pi)}$ of v1, with entropy $F(R) = h_b(R) + (1-R)\ln N$.
For $N$ above the threshold $h_b(R^\star) < R^\star \ln N$ the minimiser is the return maximiser.
So the optimality variable with a noise alphabet and our outcome $Z$ on the reward channel are one and the same random variable, and the binary case is $N = 1$.

\emph{What differs is which entropy of the run is taken.}
Given the same per-step variable, there are three different criteria.
\begin{enumerate}
\item The entropy of the long-run per-step law, $\Hmu{\occ K}$. This is our objective. It is concave in the occupancy and has a deterministic optimum.
\item The entropy of the single event "$O_t = 1$ at every step". Its probability is $\E_\pi \exp\sum_t r_t$ when $r \le 0$ is read as a log-likelihood, so this criterion is a binary entropy of an exponential utility. It is a different objective from ours.
\item The entropy of the whole symbol string, $H(Z_{1:T})$, or its rate per step. This is the criterion of PARL. It is never larger than criterion~1 and differs from it by how predictable each outcome is from the earlier ones, see the next paragraph.
\end{enumerate}
The paper uses criterion~1, and \texttt{plan.md} Section~6 records why against SMiRL and PARL.

\emph{Algorithmically the occupancy is never estimated.}
The MM loop of \eqref{eq:fp} estimates only $q$, the histogram of outcomes, and its inner step is a standard RL problem with reward $r_q(s,a) = \E_{K(\cdot \mid s,a)}[\ln q(z)]$.
The occupancy is the space on which the criterion is defined, so that \eqref{eq:fp} has a solution and that solution is a deterministic policy.
On the reward channel, write $\rho_k = q_k(\text{good symbol})$ for the current estimated success rate.
By v2 the inner reward is then $w_k\, r(s,a)$ plus a constant, with gain
\begin{equation*}
  w_k = \ln\frac{\rho_k}{1-\rho_k} + \ln N ,
\end{equation*}
which is positive exactly when $\rho_k > 1/(N+1)$.
With a policy-entropy term added, the inner step is the soft-RL problem of control as inference at inverse temperature $w_k$.
The difference from \citet{Toussaint-probabilistic06} is that their temperature is fixed and their target is always $O = 1$.
Here $w_k$ is re-estimated each iteration from the agent's own outcomes, and when $\rho_k < 1/(N+1)$ the sign of $w_k$ flips and the iteration moves towards failure.

\paragraph{Two open branches.}
\emph{Branch 1: inferring the state from outcomes.}
Suppose the regulator saw only the symbols $z_{1:t}$ and not the state.
Then the posterior over $s_t$ given $z_{1:t}$ is a hidden-Markov filter with $P$ as the transition and $K$ as the emission.
That posterior is the belief of Move~3, here over the state instead of over the disturbance.
Control as inference runs the same message passing, but conditions on the desired symbols $O_{1:T} = 1$ instead of on observed ones, so it infers how to act rather than where the agent is.
Whether the two can be combined into one inference for the error-controlled regulator is open.

\emph{Branch 2: entropy of the string against entropy of the marginal.}
By the chain rule,
\begin{equation*}
  H(Z_{1:T}) = \sum_{t=1}^{T} H(Z_t) - \sum_{t=1}^{T} I(Z_t ; Z_{<t}) .
\end{equation*}
The first sum is our criterion, up to normalisation.
The second sum is how predictable each outcome is from the earlier ones, and the string criterion gives credit for it.
Example: outcomes that alternate $A, B, A, B, \dots$ have marginal entropy $\ln 2$ and string entropy rate zero.
Our criterion charges the alternation; the string criterion does not.

Consider $Z = S$ and a deterministic policy, which both criteria may assume since both are minimised at a vertex.
The state sequence is then a Markov chain with entropy rate
\begin{equation*}
  \bar h = \sum_{s} d(s)\, H\big(P(\cdot \mid s, \pi(s))\big),
\end{equation*}
which is linear in the occupancy.
Minimising it is ordinary average-reward RL with the reward $-H(P(\cdot \mid s,a))$, a local penalty on transition noise.
Our criterion is the marginal entropy $H(d_S)$, and $H(d_S) = \bar h + I(S_{t+1} ; S_t)$.
So ours charges the transition noise plus the part of the motion that is predictable from the current state, and theirs charges the noise only.

On transients: in the long-run average neither criterion charges them.
With discounting or a finite horizon, the marginal criterion charges how spread out the transient outcomes are, while the string criterion charges only how unpredictable they are given the past, so a deterministic approach from a known start is free under the string criterion.

The internal-model reading: a periodic disturbance that the regulator can predict but not cancel produces periodic outcomes.
Their string entropy is near zero and their marginal entropy is positive.
So the string criterion rewards a regulator for having a model of the disturbance even when it cannot use it to cancel the disturbance, while the marginal criterion rewards only cancellation.
Put differently, the string criterion asks for outcomes predictable by an observer who knows the past, and ours asks for outcomes predictable by an observer who knows nothing.
Conant and Ashby's $H(Z)$ is the second.
\open{Whether the paper wants a sentence on this against PARL, or whether it is a later note.}

\section{What this gives the paper}

One regulator on $(s,b)$, drawn as \eqref{eq:level2}, replaces the two-level story: the inner regulator is $\pi$ at fixed $b$, the learner is $\pi$ curried in $b$, the model is the $B$ wire, and the occupancy is the fixed-point state of the loop.
The MM algorithm is the same loop with $B$ replaced by a point estimate and the two turns separated, which is empirical Bayes and loses $I(Z;k)$.
What both levels must model is $\psi^\pi$, which is what representation learning should keep.
The categorical notation earns its place here as a bookkeeping of wires, not as a source of theorems; the results still come from concavity on the occupancy, as Neil's note says.

\bibliographystyle{plainnat}
\IfFileExists{notes_joery/joery_survey.bib}
  {\newcommand{\surveybib}{notes_joery/joery_survey}}
  {\newcommand{\surveybib}{joery_survey}}
\bibliography{\surveybib}

\end{document}
